% ----- CHAUPAI 7 -----
\subsection*{\deva{सप्तम चौपाई} (Seventh Chaupai)}
\addcontentsline{toc}{subsection}{\deva{सप्तम चौपाई} (Seventh Chaupai) — \deva{बिद्यावान गुनी...}}

\begin{minipage}[t]{0.55\textwidth}
\vspace{0pt}

{\large \linespread{1.5}\selectfont
\deva{बिद्यावान गुनी अति चातुर।}\\
\deva{राम काज करिबे को आतुर॥}
\par}

\vspace{0.5em}

{\large \linespread{1.5}\selectfont
\telu{బిద్యావాన గునీ అతి చాతుర।}\\
\telu{రామ కాజ కరిబే కో ఆతుర॥}
\par}

\vspace{0.5em}

{\large \linespread{1.3}\selectfont
\textit{bidyāvāna gunī ati cātura |}\\
\textit{rāma kāja karibe ko ātura ||}
\par}
\end{minipage}%
\hfill
\begin{minipage}[t]{0.42\textwidth}
\vspace{0pt}

\end{minipage}

\vspace{0.5em}
\subsubsection*{\textbf{Visual Rhythm Map:}}
\begin{table}[h!]
\centering
\renewcommand{\arraystretch}{1.2}
\setlength{\tabcolsep}{0pt}
\begin{tabularx}{\textwidth}{|*{16}{>{\centering\arraybackslash\hsize=1.0\hsize}X|}}
\hline
\multicolumn{7}{|>{\centering\arraybackslash\hsize=7.0\hsize}X|}{\textbf{\guru{बि}\guru{द्या}\guru{वा}\laghu{न}} \par (7)} &
\multicolumn{3}{>{\centering\arraybackslash\hsize=3.0\hsize}X|}{\textbf{\laghu{गु}\guru{नी}} \par (3)} &
\multicolumn{2}{>{\centering\arraybackslash\hsize=2.0\hsize}X|}{\textbf{\laghu{अ}\laghu{ति}} \par (2)} &
\multicolumn{4}{>{\centering\arraybackslash\hsize=4.0\hsize}X|}{\textbf{\guru{चा}\laghu{तु}\laghu{र}} \par (4)} \\
\hline
\end{tabularx}
\vspace{-1.5em}
\end{table}

\begin{table}[h!]
\centering
\renewcommand{\arraystretch}{1.2}
\setlength{\tabcolsep}{0pt}
\begin{tabularx}{\textwidth}{|*{16}{>{\centering\arraybackslash\hsize=1.0\hsize}X|}}
\hline
\multicolumn{3}{|>{\centering\arraybackslash\hsize=3.0\hsize}X|}{\textbf{\guru{रा}\laghu{म}} \par (3)} &
\multicolumn{3}{>{\centering\arraybackslash\hsize=3.0\hsize}X|}{\textbf{\guru{का}\laghu{ज}} \par (3)} &
\multicolumn{4}{>{\centering\arraybackslash\hsize=4.0\hsize}X|}{\textbf{\laghu{क}\laghu{रि}\guru{बे}} \par (4)} &
\multicolumn{2}{>{\centering\arraybackslash\hsize=2.0\hsize}X|}{\textbf{\guru{को}} \par (2)} &
\multicolumn{4}{>{\centering\arraybackslash\hsize=4.0\hsize}X|}{\textbf{\guru{आ}\laghu{तु}\laghu{र}} \par (4)} \\
\hline
\end{tabularx}
\vspace{-1.5em}
\end{table}

\vspace{1em}
\textbf{ \deva{सरल-संस्कृतम्}:}\\
 \deva{भवान् विद्यावान्, गुणवान्, अतीव चतुरः च अस्ति।}\\
 \deva{श्रीरामस्य कार्यं कर्तुं भवान् सदैव आतुरः (उत्सुकः) अस्ति।}\\


You are highly learned, virtuous, and extremely clever/skillful.\\
You are always eager and ready to execute Lord Rama's work.


\vspace{1em}
\newpage
\textbf{\deva{प्रतिपदार्थः} (Meaning of each word):}
\begin{table}[h!]
\centering
\renewcommand{\arraystretch}{1.2}
\begin{tabularx}{\textwidth}{@{} l l X X @{}}
\toprule
\textbf{Awadhi} & \textbf{Sanskrit} & \textbf{English} & \textbf{Grammar \& Notes} \\
\midrule
\deva{बिद्यावान} \gotchaicon / \telu{బిద్యావాన}  &  \deva{विद्यावान्}  &  Learned / Scholar  & \\ 
\deva{गुनी} / \telu{గునీ}  &  \deva{गुणी (गुणवान्)}  &  Virtuous  &   \\
\deva{अति} / \telu{అతి}  &  \deva{अतीव}  &  Extremely  &   \\
\deva{चातुर} / \telu{చాతుర}  &  \deva{चतुरः}  &  Clever  &   \\
\deva{राम काज} / \telu{రామ కాజ}  &  \deva{श्रीरामस्य कार्यम्}  &  Rama's work  & Prakrit shift () \\ 
\deva{करिबे को} / \telu{కరిబే కో}  &  \deva{कर्तुम्}  &  To do (Infinitive)  & Awadhi infinitive \\
\deva{आतुर} / \telu{ఆతుర}  &  \deva{आतुरः (उत्सुकः)}  &  Eager / Ready  &   \\
\bottomrule
\end{tabularx}
\end{table}

\begin{tcolorbox}[gotchabox]
\begin{itemize}
    \item \textbf{The "Word-Initial" vs "Suffix" Rule (\deva{बिद्यावान}):} Why \deva{बिद्यावान} and not \deva{बिद्याबान}? Because the first $v$ carries the heavy front door weight of the word, it requires maximum articulatory force and hardens into the plosive $b$. The $v$ in $-vaan$ is a possession suffix (\deva{वतुँप् प्रत्यय}). Because it occurs at the end of the word where the mouth naturally relaxes (Mukha-Sukha / Lenition), it stays as the gliding breathy $v$.
\end{itemize}
\end{tcolorbox}

\begin{tcolorbox}[poetrybox]
\textbf{Praasa (Internal/External Rhyme):}\\
Line 1 focuses on his internal state (Smart / \deva{चातुर}). Line 2 focuses on his external action (Eager / \deva{आतुर}). Notice the perfect rhythm and rhyme between Chaatura and Aatura!
\end{tcolorbox}


\newpage
